\documentclass[12pt,a4paper]{article}

%% ─────────────────────────────────────────────
%%  NOTE FOR SUBMISSION:
%%  Economic Theory Bulletin requires the Springer
%%  svjour3 LaTeX class. Download it from:
%%  https://www.springernature.com/gp/authors/campaigns/latex-author-support
%%  Replace \documentclass above with:
%%  \documentclass[smallextended]{svjour3}
%%  and adapt author/institute blocks accordingly.
%%  This file compiles standalone for review purposes.
%% ─────────────────────────────────────────────

\usepackage{amsmath,amssymb,amsthm}
\usepackage{mathtools}
\usepackage{geometry}
\usepackage{xcolor}
\usepackage[colorlinks=true,linkcolor=blue,citecolor=blue,urlcolor=blue]{hyperref}
\usepackage{bm}
\usepackage{booktabs}
\usepackage{setspace}
\usepackage{enumitem}
\usepackage{microtype}
\usepackage{natbib}
\bibliographystyle{spbasic}

\geometry{margin=1.15in}
\onehalfspacing


%% ─────────────────────────────────────────────
%%  THEOREM ENVIRONMENTS
%% ─────────────────────────────────────────────
\newtheorem{theorem}{Theorem}
\newtheorem{proposition}{Proposition}
\newtheorem{corollary}{Corollary}
\newtheorem{lemma}{Lemma}
\newtheorem{remark}{Remark}
\newtheorem{definition}{Definition}

%% ─────────────────────────────────────────────
%%  CUSTOM COMMANDS
%% ─────────────────────────────────────────────
\newcommand{\E}{\mathbb{E}}
\newcommand{\Var}{\mathrm{Var}}
\newcommand{\hatpij}{\hat{p}_i^{\,j}}
\newcommand{\1}{\mathbf{1}}

%% ─────────────────────────────────────────────
\begin{document}
%% ─────────────────────────────────────────────

\title{\textbf{Strategic Marketing and GARCH Stationarity\\
under Jeffrey Updating}}

\author{Anurag Rana\\[4pt]
\small Riskcare Ltd., 25 Worship Street, London EC2A 2DX, UK\\
\small \texttt{[corresponding.author@riskcare.com]}
}

\date{\today}

\maketitle

%% ─────────────────────────────────────────────
\begin{abstract}
We study a two-player strategic investment model in which each player
markets their return process to attract external capital by shifting the
other's beliefs via Jeffrey conditioning --- without altering the true
return process. This marketing constitutes a precise form of lying: it
adds virtual up-counts to a Beta-Binomial posterior, distorting both
perceived drift and perceived volatility. The central finding concerns
the structural source of GARCH volatility dynamics. We show that the
marketing adjustment channel --- absent under pure Bayesian updating ---
provides the mean-reversion force that converts an integrated GARCH
process (IGARCH) into a stationary GARCH(1,1). Under pure Bayesian
updating the persistence coefficient equals unity and the conditional
variance follows a unit-root process; Jeffrey conditioning introduces
a finite mean-reversion coefficient driven by optimal marketing spend.
All GARCH parameters --- the ARCH coefficient, persistence, GJR
asymmetry, and baseline variance --- are derived as explicit closed-form
functions of model primitives: the binomial process parameters, the Beta
prior, the performance fee, and risk aversion. A companion paper
develops the full Stackelberg equilibrium and comparative statics.

\medskip\noindent
\textbf{JEL codes:} D82, D83, G14, C11, C58

\medskip\noindent
\textbf{Keywords:} Jeffrey conditioning, asymmetric information,
GARCH stationarity, strategic investment, Beta-Binomial beliefs,
mean reversion
\end{abstract}

\newpage

%% ─────────────────────────────────────────────
\section{Introduction}
\label{sec:intro}
%% ─────────────────────────────────────────────

GARCH models are among the most widely used tools in empirical finance,
yet their structural origins remain poorly understood. Persistence and
mean reversion in conditional variance are typically imposed
parametrically rather than derived from economic primitives. This paper
provides a microfoundation: in a two-player strategic investment model
with asymmetric information, GARCH volatility dynamics emerge from the
interaction between strategic belief manipulation, optimal capital
allocation, and a shared public history of returns.

The key finding is that the choice of belief-updating rule has
structural consequences for the nature of the GARCH process. Under
standard Bayesian updating, belief revision is driven entirely by
hard evidence --- observable returns with well-defined likelihoods.
No mechanism exists for soft evidence such as marketing or persuasion.
We show that this restriction implies a structural deficiency: the
conditional variance of the resulting portfolio follows an integrated
GARCH (IGARCH) process with unit-root persistence, which is both
empirically counterfactual and theoretically unsatisfying.

Jeffrey conditioning \citep{Jeffrey1965} allows beliefs to shift in
response to soft evidence --- evidence that changes credence without
constituting a verifiable observation. In our model, marketing spend
generates exactly this kind of shift: it adds virtual up-counts to
the Bayesian posterior without changing the true return process. This
is a lie in the precise economic sense. However, it introduces an
endogenous strategic variable whose optimal level responds to the
state of the world. This response constitutes the mean-reversion force
that restores GARCH stationarity.

The contribution of this note is to establish this Bayesian versus
Jeffrey distinction as a structural determinant of volatility dynamics
--- not a modelling detail. The result is clean: replacing Bayesian
with Jeffrey updating converts IGARCH into stationary GARCH(1,1),
with all parameters derived explicitly from economic primitives.

This complements the endogenous uncertainty literature
\citep{Kurz1994, KurzMotolese2001}, which generates GARCH from the
aggregate distribution of heterogeneous beliefs in competitive
equilibrium. Our mechanism is bilateral and strategic rather than
distributional and competitive. It also complements \citet{Kyle1985},
where informed trading converges to informational efficiency: in our
model, belief manipulation via marketing converges to truth as
accumulated history overwhelms the distortion. Both frameworks
achieve asymptotic informational efficiency through different channels.
The difference from Bayesian persuasion \citep{KamenicaGentzkow2011}
is that communication here is costly, updating is via Jeffrey rather
than Bayes, and both parties simultaneously act as senders and
receivers.

A companion paper \citep{Rana2025} develops the full model: the
Stackelberg equilibrium with closed-form solutions, comparative
statics, the monopoly limit, and the convergence theorem showing
truth dominates marketing asymptotically.

%% ─────────────────────────────────────────────
\section{Model}
\label{sec:model}
%% ─────────────────────────────────────────────

\subsection{Return Processes and Shared History}

Two players $i \in \{1,2\}$, with $j = 3-i$, each own a risky process.
Each period $t$, player $i$'s process generates a return per unit of
capital:
\begin{equation}
  r_i^t =
  \begin{cases}
    u_i & \text{with true probability } p_i,\\
    d_i & \text{with true probability } 1-p_i,
  \end{cases}
  \label{eq:returns}
\end{equation}
where $u_i > d_i$ are publicly known constants and
$p_i \in (0,1)$ is the true up-probability, private to player $i$.
Draws are i.i.d.\ across periods. True mean and variance:
\begin{equation}
  \mu_i = p_i u_i + (1-p_i)d_i,
  \qquad
  \sigma_i^2 = p_i(1-p_i)(u_i-d_i)^2.
  \label{eq:true_moments}
\end{equation}
Both are determined entirely by $p_i$. Players have different true
volatilities: $\sigma_1^2 \neq \sigma_2^2$ in general.

Both players observe the same public history of $n$ past discrete
returns from each process. Let $k_i$ be the number of up-moves in
player $i$'s history. Player $j$'s prior on $p_i$ is
$p_i \sim \mathrm{Beta}(\alpha_i^0, \beta_i^0)$. By conjugacy,
the Bayesian posterior is:
\begin{equation}
  p_i \mid \mathbf{r}_i \sim \mathrm{Beta}(\alpha_i^*, \beta_i^*),
  \qquad
  \alpha_i^* = \alpha_i^0 + k_i, \quad
  \beta_i^* = \beta_i^0 + n - k_i.
  \label{eq:posterior}
\end{equation}
Since history is public, both players agree on this posterior. Let
$S = \alpha_i^* + \beta_i^*$ denote the total pseudo-count;
$S \sim n$ for large $n$.

\subsection{Marketing as a Jeffrey Shift}

The following describes player $i$'s marketing acting on player $j$'s 
beliefs; the symmetric case --- player $j$'s marketing shifting player 
$i$'s beliefs about $p_j$ --- follows by exchanging $i$ and $j$ 
throughout. Player $i$ spends $m_i \geq 0$ on marketing. This constitutes soft
evidence in the sense of \citet{Jeffrey1965}: it shifts player $j$'s
credence over $p_i$ without being a verifiable observation and without
altering the true process $(\mu_i, \sigma_i^2)$. The Jeffrey rigidity
condition holds since $\Pr(r_i^t = u_i \mid p_i) = p_i$ is determined
by the true process, not by marketing. Formally, marketing generates a
pseudo-count shift with diminishing returns:
\begin{equation}
  \delta(m_i) = \frac{m_i}{m_i + \gamma_i}, \quad \gamma_i > 0,
  \label{eq:delta}
\end{equation}
with $\delta(0)=0$, $\delta \in (0,1)$, $\delta'> 0$,
$\delta'' < 0$, and $\lim_{m_i\to\infty}\delta(m_i) = 1$.
The Jeffrey-shifted posterior is
$\mathrm{Beta}(\alpha_i^{**}, \beta_i^{**})$ where:
\begin{equation}
  \alpha_i^{**} = \alpha_i^* + \delta(m_i), \qquad
  \beta_i^{**} = \beta_i^*.
  \label{eq:jeffrey_shift}
\end{equation}
Player $j$'s Jeffrey-updated perceived moments:
\begin{align}
  \hatpij &= \frac{\alpha_i^{**}}{\alpha_i^{**}+\beta_i^{**}}, \qquad
  \mu_i^j = \hatpij u_i + (1-\hatpij)d_i,
  \label{eq:perceived_drift}\\
  \sigma_i^{j2} &= (u_i-d_i)^2\hatpij(1-\hatpij)
  = (u_i-d_i)^2
    \frac{(\alpha_i^*+\delta(m_i))\,\beta_i^*}{(\alpha_i^*+\delta(m_i)+\beta_i^*)^2}.
  \label{eq:perceived_vol}
\end{align}

\subsection{Investment, Capital Augmentation, and Preferences}

Each player $i$ divides budget $W_i > 0$ across three uses. First,
marketing spend $m_i \geq 0$ is paid upfront, leaving an effective
investable budget $W_i^{\mathrm{eff}} = W_i - m_i$. Second, a
fraction $(1 - \alpha_i)$ of $W_i^{\mathrm{eff}}$ is deployed into
player $i$'s own process. Third, the remaining fraction
$\alpha_i \in [0,1]$ is invested cross into player $j$'s process.
The budget identity is therefore:
\begin{equation}
  W_i = m_i
        + (1-\alpha_i)\,W_i^{\mathrm{eff}}
        + \alpha_i\,W_i^{\mathrm{eff}},
  \label{eq:budget}
\end{equation}
with $W_i^{\mathrm{eff}} = W_i - m_i$ substituted throughout.

Cross-investment is productive: when player $j$ places
$\alpha_j W_j^{\mathrm{eff}}$ into player $i$'s process, that
capital is genuinely deployed and earns the same per-unit return
$r_i^t$ as player $i$'s own capital. Player $i$ retains a
performance fee $f \in (0,1)$ on the \emph{gross} return earned
by this incoming capital. Player $i$'s total effective capital
base --- own investment plus fee-adjusted incoming capital ---
is therefore:
\begin{equation}
  \tilde{A}_i = (1-\alpha_i)\,W_i^{\mathrm{eff}}
                + f\,\alpha_j\,W_j^{\mathrm{eff}}.
  \label{eq:Atilde}
\end{equation}
The first term is player $i$'s own deployment; the second is the
fee income entitlement on the capital placed by player $j$.

Conversely, when player $i$ cross-invests $\alpha_i W_i^{\mathrm{eff}}$
into player $j$'s process, player $j$ retains the fee $f$ on
that capital's gross return. Player $i$ therefore receives only
the net-of-fee fraction:
\begin{equation}
  \tilde{B}_i = \alpha_i\,W_i^{\mathrm{eff}}\,(1-f).
  \label{eq:Btilde}
\end{equation}

Terminal wealth is the sum of returns on each leg:
\begin{equation}
  W_i(T) = r_i^t\,\tilde{A}_i + r_j^t\,\tilde{B}_i
          = r_i^t\!\underbrace{\bigl((1-\alpha_i)\,W_i^{\mathrm{eff}}
              + f\,\alpha_j\,W_j^{\mathrm{eff}}\bigr)}_{\tilde{A}_i}
          + r_j^t\!\underbrace{\alpha_i\,W_i^{\mathrm{eff}}(1-f)}_{\tilde{B}_i}.
  \label{eq:wealth}
\end{equation}
The first term is the return on player $i$'s own process, earned
on both own capital and the fee entitlement from incoming capital.
The second term is the net-of-fee return on the cross-investment
in player $j$'s process.

Players maximise mean-variance utility
$\mathcal{V}_i = \E[W_i(T)] - (\rho_i/2)\Var[W_i(T)]$,
$\rho_i > 0$, using true moments $(\mu_i, \sigma_i^2)$ for their
own process and Jeffrey-shifted beliefs $(\mu_j^i, \sigma_j^{i2})$
for the other's process.

%% ─────────────────────────────────────────────
\section{Structural GARCH and the Role of Jeffrey Conditioning}
\label{sec:garch}
%% ─────────────────────────────────────────────

\subsection{Belief Dynamics and Variance Propagation}

Index periods $t = 1, 2, \ldots$ The return surprise:
\begin{equation}
  \epsilon_i^t = r_i^t - \mu_i
  = \begin{cases}
    (1-p_i)(u_i-d_i) & \text{up move, prob.\ }p_i,\\
    -p_i(u_i-d_i) & \text{down move, prob.\ }1-p_i.
  \end{cases}
  \label{eq:surprise}
\end{equation}
The posterior mean updates after each observation as:
\begin{equation}
  \Delta\hat{p}_i^{j,t}
  = \frac{\1_{\mathrm{up}}\,\beta_i^{**,t-1}
         - \1_{\mathrm{down}}\,\alpha_i^{**,t-1}}
         {S_t(S_t+1)},
  \label{eq:belief_update}
\end{equation}
with expectation-weighted approximation:
\begin{equation}
  \Delta\hat{p}_i^{j,t}
  \approx
  \frac{\bar{C}_t\,\epsilon_i^t}{(u_i-d_i)\,S_t(S_t+1)},
  \label{eq:belief_approx}
\end{equation}
where
$\bar{C}_t = p_i\beta_i^{**,t-1}/(1-p_i)
            + (1-p_i)\alpha_i^{**,t-1}/p_i$.

The belief update propagates into realised portfolio variance
$\mathcal{R}_i = \Var[W_i(T)]$ through four channels: the drift
channel, the volatility channel, the marketing adjustment channel,
and the capital augmentation channel. Let $\Lambda_t$ denote the
total sensitivity of realised variance to the current belief level.
Taking the second-order expansion of
$\E[\mathcal{R}_{i,t+1} \mid \mathcal{F}_t]$ and using
$\E[\epsilon_i^t \mid \mathcal{F}_t] = 0$ (so the linear term
vanishes), one obtains:

\begin{theorem}[Structural GARCH(1,1)]
\label{thm:garch}
Under Jeffrey conditioning with optimal marketing, the conditional
variance satisfies:
\begin{equation}
  \E[\mathcal{R}_{i,t+1} \mid \mathcal{F}_t]
  = \underbrace{\omega_i^t}_{\text{baseline}}
  + \underbrace{\alpha_i^{\mathrm{GARCH}}\,(\epsilon_i^t)^2}_{\text{ARCH}}
  + \underbrace{\beta_i^{\mathrm{GARCH}}\,\mathcal{R}_{i,t}}_{\text{persistence}},
  \label{eq:garch11}
\end{equation}
with structural coefficients derived from model primitives:
\begin{align}
  \omega_i^t &= \frac{\kappa_t}{S_t}\bar{\mathcal{R}}_i
                + \underline{\sigma}_i^2,
  \label{eq:omega}\\
  \alpha_i^{\mathrm{GARCH}}
    &= \frac{\bar{C}_t^2}{2S_t^2(S_t+1)^2}
       \frac{\partial\Lambda_t}{\partial\hat{p}_t},
  \label{eq:alpha_garch}\\
  \beta_i^{\mathrm{GARCH}}
    &= 1
       - \frac{\partial m_i^{*,t}}{\partial\hat{p}_i^{j,t}}
       \cdot
       \frac{\partial\hat{p}_i^{j,t}}{\partial\mathcal{R}_{i,t}},
  \label{eq:beta_garch}
\end{align}
where $\underline{\sigma}_i^2 = p_i(1-p_i)(u_i-d_i)^2$ is the
irreducible process variance floor, $\kappa_t$ is the strength
of the marketing mean-reversion adjustment, and $\Lambda_t$ is
derived explicitly in \citet{Rana2025}.
\end{theorem}

A GJR-GARCH extension adds an endogenous leverage term:
\begin{equation}
  \E[\mathcal{R}_{i,t+1} \mid \mathcal{F}_t]
  = \omega_i^t
  + \alpha_i^{\mathrm{GARCH}}(\epsilon_i^t)^2
  + \gamma_i^{\mathrm{GARCH}}(\epsilon_i^t)^2
    \1_{\epsilon_i^t < 0}
  + \beta_i^{\mathrm{GARCH}}\mathcal{R}_{i,t},
  \label{eq:gjr}
\end{equation}
where the sign of $\gamma_i^{\mathrm{GARCH}}$ depends on
$\hat{p}_t \gtrless 1/2$, making the leverage effect endogenous
and regime-dependent rather than parametrically imposed.

\subsection{The Central Result: Jeffrey Conditioning Restores Stationarity}

\begin{proposition}[Jeffrey vs.\ Bayesian]
\label{prop:comparison}
\begin{enumerate}[label=(\roman*)]
  \item \textbf{Under Jeffrey conditioning with optimal marketing:}
  The persistence coefficient satisfies $\beta_i^{\mathrm{GARCH}} < 1$,
  so the GARCH process \eqref{eq:garch11} is covariance-stationary
  whenever $\alpha_i^{\mathrm{GARCH}} + \beta_i^{\mathrm{GARCH}} < 1$.
  The mean-reversion force is:
  \begin{equation}
    1 - \beta_i^{\mathrm{GARCH}}
    = \frac{\partial m_i^{*,t}}{\partial\hat{p}_i^{j,t}}
      \cdot
      \frac{\partial\hat{p}_i^{j,t}}{\partial\mathcal{R}_{i,t}}
    > 0.
    \label{eq:mean_reversion}
  \end{equation}

  \item \textbf{Under pure Bayesian updating ($m_i \equiv 0$):}
  The marketing adjustment channel vanishes entirely. Therefore:
  \begin{equation}
    \beta_i^{\mathrm{GARCH}} = 1,
    \label{eq:unit_root}
  \end{equation}
  and the conditional variance follows an IGARCH process --- a
  unit-root process with no covariance stationarity.
\end{enumerate}
\end{proposition}

\begin{proof}
\textit{Part (i):} Under Jeffrey conditioning, player $i$ has a
strictly interior optimal marketing spend $m_i^* > 0$ for all
finite $n$. The derivative $\partial m_i^*/\partial\hat{p}_i^j > 0$
because higher perceived quality raises the marginal value of
attracting cross-investment (higher fee income). The derivative
$\partial\hat{p}_i^j/\partial\mathcal{R}_i \neq 0$ because higher
realised variance reduces the attractiveness of cross-investment,
reducing $\alpha_j^*$ and hence the fee income that makes marketing
profitable. The product is strictly positive, giving
$1 - \beta_i^{\mathrm{GARCH}} > 0$.

\textit{Part (ii):} Under pure Bayesian updating, soft evidence has
no mechanism to shift beliefs --- only hard signals with
well-defined likelihoods can update the posterior. Marketing
therefore cannot influence beliefs and $m_i \equiv 0$ is the only
feasible choice. With $m_i^{*,t} \equiv 0$ for all $t$, the partial
derivative $\partial m_i^{*,t}/\partial\hat{p}_i^{j,t} = 0$
identically, and equation \eqref{eq:beta_garch} reduces immediately
to $\beta_i^{\mathrm{GARCH}} = 1$.
\end{proof}

\begin{remark}
The result is structural, not parametric. Stationarity in
\eqref{eq:garch11} does not depend on the values of $\rho_i$, $f$,
$u_i$, $d_i$, or the prior parameters. It depends only on whether
the belief-updating rule permits soft evidence --- and hence optimal
marketing --- to exist as an endogenous strategic variable. Jeffrey
conditioning permits it; Bayesian updating does not.
\end{remark}

\subsection{Properties of the Structural Coefficients}

Each coefficient has a precise economic source and a clean
comparative static in $n$:

\textit{ARCH coefficient $\alpha_i^{\mathrm{GARCH}} \sim 1/n^2$:}
Longer history anchors beliefs, making individual surprises less
capable of propagating into future variance. As $n \to \infty$,
the ARCH effect vanishes.

\textit{Persistence $\beta_i^{\mathrm{GARCH}} \to 1$ as $n \to \infty$:}
Marketing spend decays at rate $1/\sqrt{n}$ as accumulated history
weakens the effectiveness of the lie (see Proposition~\ref{prop:convergence}
below). The marketing adjustment accordingly weakens, and
$\beta_i^{\mathrm{GARCH}}$ approaches unity --- though the stationarity
condition is preserved for all finite $n$ since
$\alpha_i^{\mathrm{GARCH}}$ decays faster ($1/n^2$) than
$1 - \beta_i^{\mathrm{GARCH}}$ ($1/\sqrt{n}$).

\textit{GJR asymmetry $\gamma_i^{\mathrm{GARCH}} \sim 1/n^2$:}
The endogenous leverage effect is strongest in short relationships
and disappears as history accumulates.

\textit{Baseline $\omega_i^t \geq \underline{\sigma}_i^2 > 0$:}
Bounded below by the irreducible process variance for all $n$,
preventing variance from collapsing to zero even as all dynamic
effects vanish.

\begin{corollary}[Stationarity Condition]
\label{cor:stationarity}
The process \eqref{eq:garch11} is covariance-stationary if and only if:
\begin{equation}
  \frac{\bar{C}_t^2}{2S_t^2(S_t+1)^2}
  \frac{\partial\Lambda_t}{\partial\hat{p}_t}
  <
  \frac{\partial m_i^{*,t}}{\partial\hat{p}_i^{j,t}}
  \cdot
  \frac{\partial\hat{p}_i^{j,t}}{\partial\mathcal{R}_{i,t}},
  \label{eq:stationarity}
\end{equation}
i.e., the ARCH surprise effect is dominated by the marketing
mean-reversion force. This holds whenever fee income
$f \cdot \mu_i \cdot W_j$ is sufficiently large relative to
belief sensitivity to surprises.
\end{corollary}

Table~\ref{tab:garch} summarises the regime dependence.

\begin{table}[h]
\centering
\caption{GARCH coefficient behaviour across regimes. All
parameters derived from model primitives.}
\label{tab:garch}
\begin{tabular}{@{}lcccc@{}}
\toprule
Coefficient
  & Jeffrey, finite $n$
  & Bayesian ($m_i \equiv 0$)
  & Jeffrey, $n\to\infty$ \\
\midrule
$\omega_i^t$
  & $>\underline{\sigma}_i^2$, time-varying
  & $=\underline{\sigma}_i^2$, fixed
  & $\to\underline{\sigma}_i^2$ \\
$\alpha_i^{\mathrm{GARCH}}$
  & $>0$, decays $\sim 1/n^2$
  & $>0$ (same)
  & $\to 0$ \\
$\beta_i^{\mathrm{GARCH}}$
  & $<1$ \textbf{(stationary)}
  & $=1$ \textbf{(IGARCH)}
  & $\to 1$ \\
$\gamma_i^{\mathrm{GARCH}}$
  & endogenous sign
  & endogenous sign
  & $\to 0$ \\
Stationarity
  & \textbf{Yes}
  & \textbf{No}
  & Marginal \\
\bottomrule
\end{tabular}
\end{table}

\subsection{The Lie Decays Asymptotically}

\begin{proposition}[Asymptotic Dominance of Truth]
\label{prop:convergence}
As $n \to \infty$:
\begin{equation}
  m_i^* \to 0,
  \qquad
  \hatpij \to p_i,
  \qquad
  \sigma_i^{j2} \to \sigma_i^2.
  \label{eq:convergence}
\end{equation}
Optimal marketing spend decays at rate $1/\sqrt{n}$.
Perceived beliefs and perceived volatility converge to their
true values at rate $1/n$.
\end{proposition}

\begin{proof}
The marginal belief shift satisfies
$\Pi = \partial\hatpij/\partial m_i \sim \beta_i^{**}/S^2 \sim 1/n$.
Since the marginal benefit of marketing (fee income shifted by belief
manipulation) contains $\Pi$ as a factor while marginal cost
(foregone own-process return) is bounded away from zero, the optimal
$m_i^*$ decays at rate $1/\sqrt{n}$ and reaches zero at finite
threshold $n^*$. Thereafter $\delta(m_i^*) \to 0$ and
$\hatpij \to \alpha_i^*/S \to p_i$ by the law of large numbers.
\end{proof}

\begin{remark}
This result holds regardless of the marketing budget or the
parameters of the model. Truth --- encoded in the public history of
returns --- inevitably dominates strategic distortion. The lie is
a short-run phenomenon: most effective when relationships are young
and history is short, progressively weaker as the relationship
matures, and ultimately irrelevant in the long run.
\end{remark}

%% ─────────────────────────────────────────────
\section{Discussion}
\label{sec:discussion}
%% ─────────────────────────────────────────────

\textit{Connection to RBE literature.}
\citet{Kurz1994} and \citet{KurzMotolese2001} derive GARCH from
the aggregate distribution of heterogeneous beliefs in competitive
general equilibrium. Our mechanism is bilateral and strategic rather
than distributional and competitive. Agents in RBE ignore their
effect on prices; our players explicitly account for how marketing
shifts the other's beliefs and hence allocation. The soft evidence
channel through which marketing operates is absent from RBE, which
relies on rational but diverse models of a complex economy.

\textit{Connection to Kyle (1985).}
\citet{Kyle1985} shows that informed trading produces a martingale
price with constant volatility, converging to full information
revelation by the terminal date. Our model also achieves asymptotic
informational efficiency (Proposition~\ref{prop:convergence}),
but through an opposite mechanism: Kyle's insider conceals
information and profits from price impact; our player embellishes
beliefs and profits from fee income. Neither generates stationary
GARCH --- Kyle produces constant volatility, Bayesian updating
produces IGARCH. Jeffrey conditioning is the ingredient that
restores stationarity in our setting.

\textit{Connection to Bayesian persuasion.}
\citet{KamenicaGentzkow2011} study a sender who designs signals
to persuade a Bayesian receiver. Three distinctions apply here:
belief updating is via Jeffrey rather than Bayes; communication
is costly and funded from the investment budget (making the
interior optimum for $m_i^*$ depend on economic primitives
rather than arising from signal design); and both players
simultaneously act as senders and receivers within a Stackelberg
game rather than a one-sided persuasion game.

%% ─────────────────────────────────────────────
\section{Conclusion}
\label{sec:conclusion}
%% ─────────────────────────────────────────────

We have shown that the choice between Bayesian and Jeffrey updating
has structural consequences for volatility dynamics in a strategic
investment model. Under Bayesian updating, the conditional portfolio
variance follows IGARCH with unit-root persistence. Under Jeffrey
conditioning, optimal marketing spend introduces a mean-reversion
force that converts IGARCH into stationary GARCH(1,1). All GARCH
parameters are derived as explicit functions of model primitives.
The GJR leverage asymmetry is endogenous and regime-dependent.
Truth dominates the lie asymptotically: optimal marketing decays
at rate $1/\sqrt{n}$ and beliefs converge to truth at rate $1/n$
as the shared public history accumulates.

The result identifies Jeffrey conditioning --- and the strategic
marketing it enables --- as the structural source of mean reversion
in conditional variance. This provides a microfoundation for
empirically observed GARCH stationarity that does not require
exogenous shocks to volatility or ad hoc persistence assumptions.

%% ─────────────────────────────────────────────
%%  DECLARATIONS
%% ─────────────────────────────────────────────
\paragraph{Acknowledgements.}
The author thanks colleagues at Riskcare Ltd.\ for discussions.

\paragraph{Competing Interests.}
The author has no relevant financial or non-financial interests
to disclose.

\paragraph{Funding.}
No funding was received to assist with the preparation of this
manuscript.

%% ─────────────────────────────────────────────
%%  REFERENCES
%%  ETB uses author-year citations (Name Year).
%%  Reference list: alphabetical, Springer format.
%% ─────────────────────────────────────────────
\begin{thebibliography}{10}

\bibitem[Jeffrey(1965)]{Jeffrey1965}
Jeffrey RC (1965) The logic of decision.
McGraw-Hill, New York

\bibitem[Kamenica and Gentzkow(2011)]{KamenicaGentzkow2011}
Kamenica E, Gentzkow M (2011) Bayesian persuasion.
Am Econ Rev 101(6):2590--2615.
\url{https://doi.org/10.1257/aer.101.6.2590}

\bibitem[Kurz(1994)]{Kurz1994}
Kurz M (1994) On the structure and diversity of rational beliefs.
Econ Theory 4(6):877--900.
\url{https://doi.org/10.1007/BF01213825}

\bibitem[Kurz and Motolese(2001)]{KurzMotolese2001}
Kurz M, Motolese M (2001) Endogenous uncertainty and market
volatility. Econ Theory 17(3):497--544.
\url{https://doi.org/10.1007/PL00004108}

\bibitem[Kyle(1985)]{Kyle1985}
Kyle AS (1985) Continuous auctions and insider trading.
Econometrica 53(6):1315--1335.
\url{https://doi.org/10.2307/1913210}

\bibitem[Rana(2025)]{Rana2025}
Rana A (2025) Strategic investment under asymmetric information,
Jeffrey updating, multiplicative augmentation, and endogenous
volatility. Working paper, Riskcare Ltd. Available at SSRN:
[SSRN link to be inserted]

\end{thebibliography}

\end{document}
